% SGA 3, 第 VI_B 篇, 第 1 轮英文重构.
% 范围: 引理 5.6.1.0 全文, 含底本编者注 39,
% 以及带有底本编者注 40 的过渡句.
% 底本: Polo--Gille Exp6B-13oct24.pdf, 分篇第 28--29 页
% (合订本第 369--370 页; 重版本页眉 358--359).
% 此编者新增单元止于引理 5.6.1 处的原版页边页码 353 标记之前.
% 引理 5.6.1 本身始于分篇第 29 页.
% OCR 与 Reinhold 英文 Markdown 仅为比对材料.
% 英文底本旧注: 此图为第 1 轮的 180 dpi 底本精确裁切; 下方现为原生 TikZ 图码.

\medskip
\phantomsection\label{lem:sga3-VIB-5.6.1.0}%
\noindent\textbf{引理 5.6.1.0.---}%
\footnote{\label{note:sga3-VIB-39}编者注: 我们明确写出此引理,
它用于\hyperref[lem:sga3-VIB-5.6.1]{引理 5.6.1} 的证明.}
\emph{设 \(f:X\to S\) 是局部有限呈示态射,
\(s\in S\), \(x\in X_s\). 置 \(n=\dim_x(X_s)\),
并设 \(q:\mathbb A^n_S\to S\) 是结构态射.}

\emph{假定给定一个拟有限态射 \(u:X\to\mathbb A^n_S\),
使 \(f=q\circ u\), 并假设 \(f\) 在 \(X_s\)
的一个包含 \(x\) 且维数为 \(n\) 的不可约分支
\(\mathfrak z\) 的泛点 \(z\) 处普遍开.
则 \(u\) 在点 \(x\in X\) 处普遍开.}

先注意到
\[
  (\dagger)\qquad
  \text{\emph{这些假设在每个覆盖 \(s\) 的基变换
  \(\pi:S'\to S\) 下都保持}}
\]
(即满足 \(\pi^{-1}(s)\ne\varnothing\) 的基变换).
事实上, 设 \(S'\to S\) 为这样的态射, 并设
\begin{center}
  \phantomsection\label{diag:sga3-VIB-5.6.1.0-base-change}%
  \begin{tikzcd}[row sep=large,column sep=large]
    X'
      \arrow[r,"u'"]
      \arrow[dr,"f'"'] &
    \mathbb A^n_{S'} \arrow[d,"q'"] \\
    {} &
    S'
  \end{tikzcd}
\end{center}
是基变换所得的图, \(s'\) 是 \(S'\) 中位于 \(s\) 上的点,
\(x'\) 是 \(X'_{s'}\) 中位于 \(x\) 上的点. 由于
\[
  X'_{s'}=X_s\otimes_{\kappa(s)}\kappa(s'),
\]
由泛化的提升及维数在域扩张下的不变性
(EGA~IV\(_2\), 2.3.4\textup{(i)} 和 4.1.4) 可知,
\(x'\) 包含在 \(X'_{s'}\) 的某个不可约分支 \(\mathfrak z'\) 中,
其泛点 \(z'\) 位于 \(z\) 上, 且
\[
  n\leq\dim\mathfrak z'
    \leq\dim_{x'}X'_{s'}\leq n.
\]
因此
\[
  \dim\mathfrak z'=n=\dim_{x'}X'_{s'}.
\]
由于 \(f\) 在 \(z\) 处普遍开, \(f'\) 在 \(z'\) 处普遍开.%
\footnote[1]{\label{note:sga3-VIB-5.6.1.0-generic-points}译者注:
法文底本在此句中分别印作 \(x\) 和 \(x'\).
所述假设及紧前面对泛点的选取要求此处为 \(z\) 和 \(z'\).}

置 \(Y=\mathbb A^n_S\). 由 EGA~IV\(_3\), 14.3.3.1\textup{(i)},
为证明 \(u\) 在 \(x\) 处普遍开, 只须证明对每个整数 \(r\geq0\) 及
\[
  X'=X[T_1,\ldots,T_r]
\]
中位于 \(x\) 上的每个点 \(x'\), 态射
\[
  u':X'\longrightarrow
  Y'=\mathbb A^n_S[T_1,\ldots,T_r]
\]
在 \(x'\) 处开. 现在置
\[
  S'=S[T_1,\ldots,T_r],
\]
并设 \(q'\) 为投影
\[
  Y'\simeq\mathbb A^n_{S'}\longrightarrow S'.
\]
这就是基变换 \(S'\to S\) 所得的情形.
因此由前面的观察, 归结为证明 \(u\) 在 \(x\) 处开.

置 \(y=u(x)\). 由于 \(u\) 局部有限呈示
(因为 \(f\) 和 \(q\) 都如此; 参见 EGA~IV\(_1\), 1.4.3),
EGA~IV\(_1\), 1.10.3 表明只须证明
\[
  u\bigl(\operatorname{Spec}\mathcal O_{X,x}\bigr)
  =\operatorname{Spec}\mathcal O_{Y,y}.
\]
为此可假设 \(S\) 仿射且整. 设 \(\pi:S'\to S\) 是其正规化,
并用 \(u':X'\to Y'\) 表示 \(u\) 作基变换得到的态射. 由于
\[
  \pi_Y:Y'\longrightarrow Y
\]
整且满射, 有
\[
  \operatorname{Spec}\mathcal O_{Y,y}
  =
  \bigcup_{y'}
  \pi_Y\bigl(\operatorname{Spec}\mathcal O_{Y',y'}\bigr),
\]
其中并集遍历 \(Y'\) 中位于 \(y\) 上的所有点 \(y'\).
因此只须证明, 对每个这样的 \(y'\) 及位于 \(x\) 上的每个
\(x'\in X'_{y'}\),
\[
  u'\bigl(\operatorname{Spec}\mathcal O_{X',x'}\bigr)
  =\operatorname{Spec}\mathcal O_{Y',y'}.
\]
由于这些假设在基变换 \(\pi:S'\to S\) 下保持,
遂归结为 \(S'\) 的情形; 即可假设 \(S\) 整且正规.

此时, 由 EGA~IV\(_3\), 14.3.13, 关于 \(f\) 的假设蕴含:
存在 \(X\) 的一个不可约分支 \(Z\), 它包含
\(\mathfrak z\) (因而包含 \(x\)), 支配 \(S\), 且
\[
  \dim_z(Z_s)=n=\dim(Z_\eta),
\]
其中 \(\eta\) 是 \(S\) 的泛点. 设 \(\xi\) 是 \(Z\) 的泛点.
由于 \(u\), 因而也包括
\[
  u_\eta:Z_\eta\longrightarrow\mathbb A^n_\eta,
\]
拟有限, 点
\[
  u_\eta(\xi)=u(\xi)
\]
在 \(\mathbb A^n_\eta\) 中的闭包维数为 \(n\).
因此 \(u(\xi)\) 是 \(\mathbb A^n_\eta\) 的泛点,
也就是 \(Y=\mathbb A^n_S\) 的泛点.
所以, 若 \(g\) 表示 \(u\) 在 \(Z\) 上的限制, 则
\[
  g:Z\longrightarrow Y
\]
拟有限且支配. 由于 \(Y\) 正规,
\hyperref[lem:sga3-VIB-5.6.0]{引理 5.6.0} 表明 \(g\) 开.
故 \(u\) 在 \(Z\) 的每个点处开, 特别地在 \(x\) 处开. 因此,
\[
  u\bigl(\operatorname{Spec}\mathcal O_{X,x}\bigr)
  =\operatorname{Spec}\mathcal O_{Y,y},
\]
这就完成了\hyperref[lem:sga3-VIB-5.6.1.0]{引理 5.6.1.0} 的证明.

下面的引理可更有效地替代 EGA~IV\(_3\), 14.5.10,%
\footnote{\label{note:sga3-VIB-40}编者注: 原文标为 19.5.10, 已予改正.}
因为它不依赖于 Noether 假设.
